% Introduction — see issue #21.
\section{Introduction}
\label{sec:introduction}

Alzheimer's disease (AD) is a progressive neurodegenerative disorder that affects millions of people worldwide and places a substantial burden on healthcare systems \parencite{taiyebkhosroshahi2025explainable}. It has no cure, but treatment started in the early stage can slow its progression, which makes early diagnosis a clinical priority \parencite{nibila2026artificial}. Structural magnetic resonance imaging (MRI) is a widely used, non-invasive way to observe the brain atrophy that accompanies the disease \parencite{frisoni2010clinical}, and deep learning, in particular convolutional neural networks (CNNs), has become the dominant approach to classify these scans automatically as cognitively normal, mild cognitive impairment or AD \parencite{nibila2026artificial}.

Despite their reported accuracy, these models are hard to adopt in clinical practice. A deep network behaves as a black box: it returns a label without showing which parts of the scan led to it, and clinicians have little reason to trust a prediction they cannot inspect \parencite{taiyebkhosroshahi2025explainable}. Reported performance is also difficult to compare and sometimes inflated, since many studies split their data by image instead of by subject, so that scans of the same brain appear in both training and test sets \parencite{wen2020convolutional}. Explainable artificial intelligence (XAI) techniques address the first problem, but their use in AD neuroimaging remains limited by small datasets and the lack of standardization \parencite{taiyebkhosroshahi2025explainable}.

This work classifies AD from structural MRI with a CNN whose predictions can be interpreted, evaluated with a subject-level split to avoid data leakage. Our contributions are (i) a review of the recent literature on AD classification, summarized in a comparative matrix; (ii) an interpretable CNN pipeline for the CN, MCI and AD classes; and (iii) an assessment of both its performance and of the image regions that support its decisions. The rest of the paper is organized as follows. \Cref{sec:theoretical_framework} presents the theoretical background, \cref{sec:related_work} reviews related work, \cref{sec:methods} describes the methodology, \cref{sec:results} reports the results, \cref{sec:discussion} discusses them, and \cref{sec:conclusion} concludes.
